% =============================================================================
%  HANDOVER FOR v4 -- the Chapter 7 (Discussion) prose v3 wrote, v3.0-v3.6
% -----------------------------------------------------------------------------
%  Removed from v3's chapters/07_discussion.tex in v3.7 (2026-09-14): the author
%  moved the Discussion to v4 (Working_Space/DRAFT_OWNERSHIP.md). Kept here, in
%  v3's own folder, so v4 can mine it without v3 writing into v4/.
%
%  NOT READY TO PASTE. It predates the naming rules of v2.15 / v3.7: it still says
%  "arm", "the consistency target", "the funnel", "regime"/"all-pass", and it
%  describes the ORIGINAL corridor scene, which corridor_v2 (Gen15 U16,
%  2026-09-14) has replaced. Translate with
%  Writing/Auxiliary/Naming/TRANSLATION_20260914_dev_jargon_to_scientific.md
%  and re-check every number against the current results before reuse.
% =============================================================================

% =============================================================================
%  CHAPTER 7 -- DISCUSSION                                             [v3 OWNS]
% -----------------------------------------------------------------------------
%  All four sections are writable today (DATASTATUS section 7). Two of them --
%  sec:disc:negative and sec:disc:threats -- carry material that exists ONLY
%  because results were kept rather than discarded, and they are the reason the
%  contribution list in Chapter 1 says two of six contributions are negative or
%  methodological.
% =============================================================================

\chapter{Discussion}\label{ch:discussion}

\section{Interpretation}\label{sec:disc:interpretation}

\subsection*{Two goals of this thesis pull in opposite directions}

The engine result says: push the step budget down, because that is where the compute factor lives
--- about fifteen of the twenty-one-fold advantage in \autoref{sec:res:state} comes from running at
a budget the baseline cannot reach, and only about $1.4\times$ from the engine at equal budget. The
constraint result says the opposite: the in-ODE endpoint arm needs \emph{genuine} steps to exist at
all, and below the threshold of \autoref{sec:res:constraints:degenerate} it silently becomes
sample-then-project. One goal wants the budget at one; the other requires it above a floor.

They meet at a moderate budget with a low activation threshold. That is not a compromise imposed by
the data --- it follows from the two mechanisms. The generator's quality is flat in the budget, so
buying additional steps costs compute and returns nothing on quality; what those steps buy is
\emph{opportunities to project}, which is a property of the constraint arm, not of the engine.
Choosing the budget is therefore a constraint-enforcement decision that happens to be spelled as a
sampler parameter, and this is the sharpest practical consequence of the whole thesis: with a
diffusion engine the two decisions are welded together at training time, and with a transport engine
they come apart.

\subsection*{Where the advantage actually comes from}

A single headline number fuses three claims. \emph{Transport
instead of diffusion} buys about $1.4\times$ at equal budget and an exact tie on plan quality --- a
real but modest engine effect, consistent with two networks of identical size doing comparable work
per evaluation. \emph{The few-step objective} is what separates on quality, and only where the task
has the dynamic range to show separation at all. \emph{The budget ladder} is where the order of
magnitude is. Reporting only the product of the three, as a single ``$31\times$'', would be true and
would explain nothing.

\subsection*{Projection is not universally worth its price}

The aerial entry shows the cleanest version of a result that is implicit in all three: whether the
projector earns its cost is an engine-level question, not a global one. On that entry the
consistency target is at zero without projection --- it does not function without it ---
MeanFlow is marginal, and plain flow matching's best configuration is \emph{unprojected}. A
generator whose plans are nearly feasible pays the projector's price for very little, and one whose
plans are not is rescued by it. This is the same quantity that appears in
\autoref{sec:res:constraints} as the reason the endpoint arm is cheap --- how far the point handed
to the solver is from the feasible set --- seen from the other end.

\section{Negative and Inconclusive Results}\label{sec:disc:negative}

\subsection*{The consistency target does not improve on MeanFlow}

This was a planned rung of the engine ladder and it did not hold. Stating it precisely matters more
than stating it briefly, because the imprecise versions are all wrong in one direction or the other.

\textbf{What was measured.} In four environments: on the state-based benchmark it is
\emph{indistinguishable} from MeanFlow at low budget --- nominally ahead, $p = 0.231$, and a wash
over all six projection rules; on the vision-conditioned entry it ties at a budget of two and is
\emph{worse} at twenty; on the aerial pillars scene it finishes \emph{behind plain flow matching},
the only engine with no configuration above $0.8$; and on the corridor scene nothing separates at
all. \textbf{Better nowhere; equal where the task cannot separate.}

\textbf{Name the floor.} ``The consistency target'' is not one arm. Two floors were trained,
$\alpha = 0.2$ and $\alpha = 0.05$, and they behave differently: on the vision-conditioned entry the
significant result at a budget of twenty ($p = 0.0215$) is the $\alpha = 0.05$ arm, while
$\alpha = 0.2$ shows a mean deficit of $3.4\times$ that the sign test does \emph{not} reject. A
sentence that says ``$\alpha$-Flow is worse'' without naming the floor is not reproducible.

\textbf{The mechanism.} The target closes its finite difference with a stop-gradient copy of the
network, so it inherits a fraction of that network's own error, at a rate set by $\alpha$. The
inherited term does not shrink as the interval shrinks, so it is independent of the step budget, and
the objective admits self-consistent fields that are wrong. That predicts what
\autoref{sec:res:visual} measures: MeanFlow improves from $0.2354$\,m to
$0.0741$\,m as the budget rises tenfold, while both consistency arms stay flat or worsen. An
engine whose error cannot be integrated away is not one that a larger budget rescues.

\guard{One selection effect must travel with the low-budget table:} two floors were trained and the
better one is quoted there, on a single seed. Both arms are reported, and the comparison is a pair
rather than an ordering.

\subsection*{The quality version of the flow-matching claim is withdrawn}

An earlier form of the engine ladder had plain flow matching beating the diffusion baseline on
quality. That does not survive: on the vision-conditioned entry the two are $0.0055$\,m apart inside
a band of roughly $0.4$\,m at $p = 1.0000$, an exact tie, and at matched budget on the state-based
benchmark the pooled paired difference over thirty pairs is $-0.012$. What survives is the
\emph{cost} form of the claim, on the state-based benchmark, and this thesis states only that form.
The two have been conflated in this area before, including in earlier drafts of this work.

\subsection*{A scene retired as a ranking instrument, and kept as something better}

The aerial s-curve scene was built to be the hard case and cannot rank engines: its ceiling is
$0.100$ across forty configurations, no arm clears the second funnel stage, and an ordering
previously reported on it was \emph{tracker-confounded and is retracted}. Raising the budget beyond
five is non-monotone and does not help. It is kept as the demonstration of where the control stack runs out: the retraction is the evidence that the aerial entry's limit is the tracker and not the planner.

\subsection*{Two results that carry no claim}

A safety advantage for endpoint projection at the largest budget on the vision-conditioned entry
would rest on ten of ten against nine of ten zero-violation rollouts --- \emph{one} discordant
rollout, at $p = 1.0000$. The improved-MeanFlow variant explored early in this project is out of
scope rather than refuted: it was never carried into the matched-architecture protocol, so there is
no comparable evidence about it either way.

\section{Threats to Validity}\label{sec:disc:threats}

\subsection*{Task saturation, found in this work's own results}

The most transferable methodological finding here is uncomfortable and general.
\textbf{A benchmark can be simultaneously feasible, well-instrumented, and completely unable to
support the comparison being run on it}, and nothing in the resulting numbers advertises the fact.

The aerial corridor scene is the worked example. Every arm passes, projected and unprojected alike;
violations read exactly zero on every row while the very same counters report violations in the
hundreds on a sibling scene, so the instrument is demonstrably working. The primary metric there
silently degenerates to plain success, and the engine ordering it produces --- a spread of 2.4
control steps against a standard deviation near 5.0 --- is noise presented in the format of a result.
Worse, the projector on that scene spends 155 to 202\,ms per control step removing violations that
never existed. A reader shown only that table would see a well-behaved constrained-control
experiment.

The general form: \textbf{an engine comparison on a saturated benchmark measures the benchmark, not
the engine}, and reporting the regime is the cheapest available defence. It costs one column --- the
violation count on the \emph{unprojected} arm --- and it is why \autoref{sec:setup:tasks} defines
the regime taxonomy before any result is shown. The same defence would have caught the low-budget
ordering on the state-based benchmark, which shows the same shape independently.

\subsection*{Downstream metrics can hide the generator}

Every number in \autoref{ch:results} is measured after a projector and a tracking controller have
acted on the plan. Both are low-pass: the projector pulls a plan toward the feasible set, and the
controller converts a sequence of setpoints into commands through inverse kinematics or a cascaded
geometric law. Either can absorb a great deal of incoherence in the plan it is given.

\autoref{fig:raw-plans} shows what that conceals. At one network evaluation the diffusion baseline's
plans are a high-frequency scribble across the workspace, and its \emph{commanded} traces are
nonetheless smooth --- the controller integrates the noise away. A reader given only the executed
trajectories and the success rates would see two comparable planners. Turning the projector off
separates them by a factor of $4.7$ in violations at matched compute.

The general form: \textbf{an evaluation that measures only what survives the control stack measures
the control stack's tolerance as much as the planner's quality.} Reporting the unprojected arm costs
one extra configuration per cell and is the cheapest available defence --- it is why the funnel of
\autoref{sec:setup:metrics} evaluates its first stage there, and why this thesis reports raw plans at
the budget where the engines are supposed to differ. The same argument bounds what the aerial entry
can show: with a tracking error of 0.30--0.49\,m, that stack absorbs more than any of the scenes
leave as clearance.

\subsection*{A guided sampler that is not guided}

The condition of \autoref{sec:res:constraints:degenerate} is a threat to validity of a whole class of
published comparisons, not only of ours. Where the number of genuine steps falls to zero, an in-loop
guided sampler computes sample-then-project while being reported as guidance --- and the failure is
silent, because both arms run, both produce plans, and both report timings. This work's
instrumentation prints the count at run time and every configuration in \autoref{ch:results} is
tagged by it. The condition is ours and is not claimed for the source method, whose published
configuration sits comfortably above the floor.

\subsection*{Scope and coverage}

\emph{One entry carries the multi-seed evidence.} The state-based benchmark is the only entry with
five seeds; the other two are single-seed, one of them compensating with pairing and one not
compensating at all. Every claim outside the first entry is qualified in place, and the aerial
ranking is explicitly not banked.

\emph{The baseline's coverage is asymmetric.} On the hardest geometry of the first entry the
baseline has one seed against the transport engines' five. The geometry-first aggregation makes the
comparison possible; it does not make the baseline's hardest cell well-estimated.

\emph{Two increments were added in one step.} The second entry changes both the control problem and
the observation modality relative to the first. Its results therefore support ``the result survives
a harder, differently-observed task'' and \emph{not} ``the result transfers from state to vision'',
because the state version of that task was never run. \autoref{sec:disc:limitations} states what
would separate them.

\emph{Timing is measured on a shared, CPU-limited node.} The reasons are given with the machine in
\autoref{sec:setup:protocol}. Ratios are same-node and same-condition, and the load-bearing ones are
factors of twenty and thirty; the figures for which this caveat could matter are the near-unity ones,
and no claim is built on those.

\emph{One solver failure is open.} The constraint program fails to converge at one specific transport
time in every configuration of the endpoint arm on the vision-conditioned entry --- twenty-four of
twenty-four items across three objectives. Engine-side causes are therefore ruled out; the remaining
suspect is the conditioning of the constraint Jacobian. It is documented rather than resolved.

\section{Limitations}\label{sec:disc:limitations}

\textbf{The control stack has an edge, and it was found.} The aerial s-curve scene is not trackable
by the deployed controller at any budget tested. The binding element is the tracker, not the
planner: that is what the retraction of the earlier ordering on that scene demonstrates. The
consequence for this thesis is bounded and specific --- the embodiment transfer of \textbf{RQ4}
holds on the scenes the control stack can hold, and stops where the reference becomes too aggressive
for it.

\textbf{One entry sits in the all-fail regime, and retraining would not move it.} No arm satisfies
the constraint set on the vision-conditioned entry. Three of the four candidate causes are
insensitive to training: an over-strict constraint set is a property of the geometry, the
reproducibility floor is run-to-run variance, and ten contexts is a protocol choice --- all three
unchanged by new weights. Only checkpoint quality is a training question, and it should be tested
last and with a stated hypothesis, because a retrain that also changes the constraint set confounds
the two effects and answers neither. \guard{If a retrain does not move the entry out of the all-fail
regime, that is itself the finding: the over-strictness is a property of the constraint set, the
funnel is the correct treatment.}

\textbf{The state-based leg of the second entry does not exist.} The three-dimensional state version
of the aligning task has an environment but no configuration, no checkpoints and no evaluation
corpus. It is what would separate \emph{harder control} from \emph{harder perception}, and without
it the two increments are confounded. This is the largest single gap between the evidence and the
claim structure.

\textbf{The aerial ranking is one seed of one scene.} Two to three additional seeds would settle it.
Until they land, the aerial entry contributes the embodiment argument, the regime taxonomy, and the
projection cost and safety result --- all of which are independent of the ranking --- and the
ranking itself carries its seed depth in every sentence.

\textbf{No quantitative smoothness metric exists in this pipeline.} Plan quality is measured by
control steps to completion, by constraint satisfaction, and --- on the unprojected arm --- by
violation counts and by inspection of the plans themselves (\autoref{fig:raw-plans}). Two plans with
the same step count and the same violation record can still differ in curvature, in jerk and in the
control effort they demand, and none of those is measured. The specific missing measurement is jerk,
path length or curvature computed over the saved plan files, which requires no new experiment ---
only a pass over plans already on disk. Until it is taken, the plan-quality argument of
\autoref{sec:res:state} rests on violation counts, which are quantitative, and on
\autoref{fig:raw-plans}, which is not.
